\documentclass{labreport}

\usepackage{microtype}
\usepackage{enumitem}
\usepackage{pifont}
\usepackage{tabularx}
\usepackage{booktabs}

% Headheight adjustment for fancyhdr
\setlength{\headheight}{28pt}
\addtolength{\topmargin}{-13pt}

% Colors
\definecolor{codebg}{RGB}{248, 249, 251}
\definecolor{codeborder}{RGB}{218, 222, 229}
\definecolor{linkcolor}{RGB}{20, 60, 120}
\definecolor{zshgreen}{RGB}{0, 160, 80}
\definecolor{zshred}{RGB}{200, 30, 30}

% Hyperref setup with subtle link colors
\hypersetup{
    colorlinks=true,
    linkcolor=linkcolor,
    filecolor=linkcolor,
    urlcolor=linkcolor,
    citecolor=linkcolor
}

% Code and terminal listings
\lstdefinestyle{bash}{
    basicstyle=\ttfamily\footnotesize,
    backgroundcolor=\color{codebg},
    frame=single,
    framerule=0.5pt,
    rulecolor=\color{codeborder},
    breaklines=true,
    breakatwhitespace=false,
    showstringspaces=false,
    keepspaces=true,
    columns=flexible,
    numbers=none,
    xleftmargin=3pt,
    xrightmargin=3pt,
    framesep=4pt
}

\lstdefinestyle{c}{
    language=C,
    basicstyle=\ttfamily\footnotesize,
    backgroundcolor=\color{codebg},
    frame=single,
    framerule=0.5pt,
    rulecolor=\color{codeborder},
    keywordstyle=\color{blue!85!black}\bfseries,
    commentstyle=\color{gray!80!black}\itshape,
    stringstyle=\color{teal!80!black},
    numbers=left,
    numberstyle=\tiny\color{gray!80!black},
    numbersep=8pt,
    breaklines=true,
    breakatwhitespace=false,
    showstringspaces=false,
    keepspaces=true,
    columns=flexible,
    tabsize=2,
    xleftmargin=14pt,
    xrightmargin=3pt,
    framesep=4pt
}

\lstset{
    extendedchars=true,
    literate=
        {➜}{{\textcolor{zshgreen}{\ding{220}}}}{1}
        {├─}{|-\,}{2}
        {└─}{\`-\,}{2}
        {├}{|}{1}
        {─}{-}{1}
        {└}{\`}{1}
        {✗}{{\textcolor{zshred}{\ding{55}}}}{1}
}

% Section styling for clean document structure
\titleformat{\section}
  {\large\bfseries}{\thesection}{0.8em}{}[\vspace{1pt}\hrule height 0.5pt]
\titleformat{\subsection}
  {\normalsize\bfseries}{\thesubsection}{0.8em}{}
\titleformat{\subsubsection}
  {\small\bfseries}{\thesubsubsection}{0.8em}{}

\begin{document}

\labtitle
    {8}
    {The Same Calls Inside Your Own OS: Process Management and IPC in xv6}
    {st. gr. FAF-243}
    {Șoimu Ionuț}
    {asist. univ.\\[0.15cm] Reițman Patricia}

\tableofcontents
\newpage

% ==============================================================================
\section{Lab 8: Track B -- The Same Calls Inside Your Own OS (xv6)}

\subsection{Before you start}

\subsubsection*{Purpose and Scope}
In Laboratory Work 7, we investigated process creation, binary image replacement, and stream redirection on a host Linux system using the GNU C Library (\texttt{glibc}) and the Linux 6.x kernel. The primary objective of Laboratory Work 8 is to bridge user-space software abstractions with low-level kernel internals by implementing and evaluating the \textbf{exact same process management and Inter-Process Communication (IPC) primitives} inside \textbf{xv6-riscv}, a modern re-implementation of Dennis Ritchie and Ken Thompson's Unix Version 6 targeting 64-bit RISC-V hardware.

Unlike Linux user-space programs that rely on thousands of pages of standard library scaffolding, an xv6 user program executes directly on top of the bare-metal kernel interface. There is no \texttt{glibc}, no dynamic linker, and no complex runtime: every invocation of \texttt{fork()}, \texttt{exec()}, \texttt{wait()}, \texttt{pipe()}, \texttt{read()}, and \texttt{write()} transitions via RISC-V \texttt{ecall} instructions directly into the xv6 kernel trap handler.

\subsubsection*{xv6 Workspace and Build Integration}
Work is conducted in the repository tree \texttt{\textasciitilde/xv6-riscv}. In xv6, user-space utilities are compiled as standalone ELF binaries from \texttt{user/NAME.c} and embedded into the initial file system image (\texttt{fs.img}) created by \texttt{mkfs}. Adding a new program requires inserting \texttt{\$U/\_NAME\textbackslash} into the \texttt{UPROGS} definition within the \texttt{Makefile}:

\begin{lstlisting}[style=bash]
UPROGS=\
    $U/_cat\
    $U/_echo\
    $U/_forktest\
    $U/_grep\
    $U/_init\
    $U/_kill\
    $U/_ln\
    $U/_ls\
    $U/_mkdir\
    $U/_rm\
    $U/_sh\
    $U/_stressfs\
    $U/_usertests\
    $U/_grind\
    $U/_wc\
    $U/_zombie\
    $U/_run\
    $U/_pingpong\
    $U/_primes\
\end{lstlisting}

Every target line begins with a leading tab character and concludes with a line-continuation backslash (\texttt{\textbackslash}). We verified the compilation and execution of xv6 using QEMU:

\begin{lstlisting}[style=bash]
➜  xv6-riscv make qemu
qemu-system-riscv64 -machine virt -bios none -kernel kernel/kernel -m 128M -smp 3 -nographic -global virtio-mmio.force-legacy=false -drive file=fs.img,if=none,format=raw,id=x0 -device virtio-blk-device,drive=x0,bus=virtio-mmio-bus.0

xv6 kernel is booting

hart 1 starting
hart 2 starting
init: starting sh
$ 
\end{lstlisting}

% ==============================================================================
\subsection{Part 1. run, inside xv6}

\subsubsection{Concept and Program Structure}
In Unix-like systems, the command-line shell executes commands through a classic three-step sequence:
\begin{enumerate}
    \item \textbf{Process Spawning (\texttt{fork}):} The shell duplicates its own address space, producing an identical child process.
    \item \textbf{Program Invocation (\texttt{exec}):} The child process executes \texttt{exec()}, replacing its virtual memory image with the targeted binary.
    \item \textbf{Process Synchronization (\texttt{wait}):} The parent suspends execution until the child completes, retrieving its termination status.
\end{enumerate}

In xv6, these steps are executed through the minimal headers \texttt{kernel/types.h} and \texttt{user/user.h}. We implemented \texttt{user/run.c}:

\begin{lstlisting}[style=c]
#include "kernel/types.h"
#include "user/user.h"

// the three steps of the shell, inside xv6:  echo hello from the child
int
main(int argc, char *argv[])
{
  printf("parent: my pid is %d\n", getpid());
  int pid = fork();                     // step 1: copy myself
  if(pid < 0){
    printf("fork failed\n");
    exit(1);
  }
  if(pid == 0){                         // only the child runs this
    printf("child: my pid is %d\n", getpid());
    char *args[] = { "echo", "hello", "from", "the", "child", 0 };
    exec("echo", args);                 // step 2: become echo
    printf("exec failed\n");            // runs only if exec failed
    exit(1);
  }
  int status;
  wait(&status);                        // step 3: wait for the child
  printf("parent: child %d ended with status %d\n", pid, status);
  exit(0);
}
\end{lstlisting}

\subsubsection{Execution and Experiments}

\begin{enumerate}[leftmargin=*]
    \item \textbf{Default Execution with \texttt{echo}:}
\begin{lstlisting}[style=bash]
$ run
parent: my pid is 3
child: my pid is 4
hello from the child
parent: child 4 ended with status 0
\end{lstlisting}
    The parent process (PID 3) invokes \texttt{fork()}, allocating a newly initialized \texttt{proc} struct with PID 4. The child prints its PID, transforms into \texttt{echo} via \texttt{exec("echo", args)}, and outputs the argument string. When \texttt{echo} finishes, the parent's \texttt{wait(\&status)} wakes up, reporting exit status 0.

    \item \textbf{Modifying the Child to Execute \texttt{ls}:}
    We modified the argument array and \texttt{exec} invocation:
\begin{lstlisting}[style=c]
    char *args[] = { "ls", 0 };
    exec("ls", args);                 // step 2: become ls
\end{lstlisting}
    Rebuilding and launching \texttt{run} in QEMU produced:
\begin{lstlisting}[style=bash]
$ run
parent: my pid is 3
child: my pid is 4
.              1 1 1024
..             1 1 1024
README         2 2 2441
cat            2 3 35864
echo           2 4 34760
sleep          2 5 34440
args           2 6 34424
mystat         2 7 35616
find           2 8 42048
run            2 9 35024
pingpong       2 10 35800
primes         2 11 36464
forktest       2 12 17704
grep           2 13 43288
init           2 14 35112
kill           2 15 34664
ln             2 16 34456
ls             2 17 41944
mkdir          2 18 34712
rm             2 19 34696
sh             2 20 57072
stressfs       2 21 35536
usertests      2 22 201520
grind          2 23 50336
wc             2 24 36832
zombie         2 25 34008
logstress      2 26 36680
forphan        2 27 35472
dorphan        2 28 34936
sync           2 29 34112
console        3 30 0
parent: child 4 ended with status 0
\end{lstlisting}

    \item \textbf{Executing a Non-Existent Binary (\texttt{nosuchprog}):}
    We tested error recovery by replacing the executable target with an invalid name:
\begin{lstlisting}[style=c]
    char *args[] = { "nosuchprog", 0 };
    exec("nosuchprog", args);
\end{lstlisting}
    Executing in QEMU:
\begin{lstlisting}[style=bash]
$ run
parent: my pid is 3
child: my pid is 4
exec failed
parent: child 4 ended with status 1
\end{lstlisting}
\end{enumerate}

\subsubsection*{Observe:}
\begin{enumerate}[leftmargin=*]
    \item \textbf{Write the three calls that run.c uses in xv6 and the three calls run.c used on Linux in Lab 7. What is the same?}\\[0.2em]
    \textbf{xv6 calls:}
    \begin{itemize}[noitemsep]
        \item \texttt{fork()}
        \item \texttt{exec(path, argv)}
        \item \texttt{wait(\&status)}
    \end{itemize}
    \textbf{Linux calls (Lab 7):}
    \begin{itemize}[noitemsep]
        \item \texttt{fork()}
        \item \texttt{execlp("ls", "ls", "-l", (char *)NULL)}
        \item \texttt{waitpid(pid, \&status, 0)} (or \texttt{wait(\&status)})
    \end{itemize}
    \textbf{What is the same:}\\
    The high-level process execution model and interface are fundamentally identical. In both operating systems:
    \begin{itemize}[noitemsep]
        \item \texttt{fork()} separates creation from execution: it clones the parent's address space, returning 0 in the child and the child's PID in the parent.
        \item \texttt{exec()} loads an ELF binary from storage, discards the old user memory pages, initializes new page tables, and begins execution at the entry point without altering the PID.
        \item \texttt{wait()} puts the parent to sleep until a child terminates, recovering the child's exit status and allowing the kernel to reclaim its process table entry.
    \end{itemize}

    \item \textbf{Who prints "hello from the child": the code of run, or echo?}\\[0.2em]
    The program \textbf{\texttt{echo}} prints the message. Once \texttt{exec("echo", args)} succeeds, the code, data, and stack segments of \texttt{run.c} in the child process no longer exist in memory. The CPU registers and program counter are pointed to the entry address of the \texttt{echo} binary. The printed string is produced by the \texttt{printf()} call inside \texttt{echo.c}'s \texttt{main()} routine.

    \item \textbf{Why is the status 1 here and 127 on Linux? (Look at the exit line after exec in both programs.)}\\[0.2em]
    The exit code received by the parent process is determined entirely by the integer parameter passed to \texttt{exit()} by the child when \texttt{exec} fails:
    \begin{itemize}
        \item In xv6 (\texttt{user/run.c}):
\begin{lstlisting}[style=c]
    printf("exec failed\n");
    exit(1);
\end{lstlisting}
        The child explicitly terminates with argument \texttt{1}. The xv6 kernel stores this value in \texttt{p->xstate}, and \texttt{wait(\&status)} copies \texttt{p->xstate} into the parent's status variable.
        \item In Linux (\texttt{run.c}):
\begin{lstlisting}[style=c]
    perror("exec");
    exit(127);
\end{lstlisting}
        The child explicitly terminates with argument \texttt{127}, following the standard POSIX shell convention where exit code 127 indicates that a specified command or binary was not found.
    \end{itemize}
\end{enumerate}

% ==============================================================================
\subsection{Part 2. pingpong}

\subsubsection{Inter-Process Communication Across Unidirectional Pipes}
A pipe is an anonymous byte stream provided by the operating system kernel. In xv6, calling \texttt{pipe(p)} allocates a circular buffer in kernel memory and returns two file descriptors:
\begin{itemize}[noitemsep]
    \item \texttt{p[0]}: the read descriptor.
    \item \texttt{p[1]}: the write descriptor.
\end{itemize}

Because pipes in xv6 are strictly unidirectional, two distinct pipes are necessary to establish bidirectional communication between parent and child without race conditions:
\begin{enumerate}[noitemsep]
    \item \texttt{p2c} (Parent-to-Child): parent writes to \texttt{p2c[1]}; child reads from \texttt{p2c[0]}.
    \item \texttt{c2p} (Child-to-Parent): child writes to \texttt{c2p[1]}; parent reads from \texttt{c2p[0]}.
\end{enumerate}

\subsubsection{Implementation of user/pingpong.c}
We implemented \texttt{user/pingpong.c}:

\begin{lstlisting}[style=c]
#include "kernel/types.h"
#include "user/user.h"

int
main(int argc, char *argv[])
{
  int p2c[2], c2p[2];     // parent-to-child, child-to-parent
  char b = 'x';

  if(pipe(p2c) < 0 || pipe(c2p) < 0){
    fprintf(2, "pipe failed\n");
    exit(1);
  }

  int pid = fork();
  if(pid < 0){
    fprintf(2, "fork failed\n");
    exit(1);
  }

  if(pid == 0){
    close(p2c[1]); close(c2p[0]);       // the ends the child does not use
    if(read(p2c[0], &b, 1) != 1){
      fprintf(2, "child read failed\n");
      exit(1);
    }
    printf("%d: received ping\n", getpid());
    if(write(c2p[1], &b, 1) != 1){
      fprintf(2, "child write failed\n");
      exit(1);
    }
    close(p2c[0]);
    close(c2p[1]);
    exit(0);
  }

  close(p2c[0]); close(c2p[1]);         // the ends the parent does not use
  if(write(p2c[1], &b, 1) != 1){
    fprintf(2, "parent write failed\n");
    exit(1);
  }
  if(read(c2p[0], &b, 1) != 1){
    fprintf(2, "parent read failed\n");
    exit(1);
  }
  printf("%d: received pong\n", getpid());
  close(p2c[1]);
  close(c2p[0]);
  wait(0);
  exit(0);
}
\end{lstlisting}

\subsubsection{Execution in xv6}
We compiled xv6 with \texttt{pingpong} included in \texttt{UPROGS} and tested it in QEMU:

\begin{lstlisting}[style=bash]
$ pingpong
6: received ping
5: received pong
$
\end{lstlisting}
The output order confirms proper synchronization: the child (PID 6) prints \texttt{received ping} first, followed by the parent (PID 5) printing \texttt{received pong}.

\subsubsection*{Observe:}
\begin{enumerate}[leftmargin=*]
    \item \textbf{Why do you need TWO pipes, and not one for both directions?}\\[0.2em]
    A Unix/xv6 pipe is a single unidirectional FIFO queue. If a single pipe \texttt{p[2]} were shared for both directions:
    \begin{itemize}[noitemsep]
        \item When the parent writes a byte into \texttt{p[1]} and immediately calls \texttt{read(p[0])}, the parent could read back its own byte before the child process is scheduled.
        \item If both processes read simultaneously from the same pipe buffer, non-deterministic race conditions occur.
    \end{itemize}
    Allocating two separate pipes establishes two disjoint kernel buffers, isolating traffic from parent-to-child (\texttt{p2c}) from traffic in the return direction (\texttt{c2p}).

    \item \textbf{Nothing in your code says "wait until the child has replied". Which call makes the parent wait?}\\[0.2em]
    The blocking system call \textbf{\texttt{read(c2p[0], \&b, 1)}} forces the parent to wait. In xv6, reading from an empty pipe puts the process into the \texttt{SLEEPING} state on the pipe's internal condition channel (\texttt{piperead()}). The parent cannot proceed until the child writes the byte into \texttt{c2p[1]}, triggering a kernel \texttt{wakeup()} on the channel.
\end{enumerate}

% ==============================================================================
\subsection{Optional Extension: primes (Concurrent Prime Sieve)}

\subsubsection{Algorithm and Process Architecture}
The Prime Sieve pipeline (originally conceived by Doug McIlroy, inventor of Unix pipes) generates prime numbers using concurrent processes arranged as a pipeline:
\begin{itemize}[noitemsep]
    \item The root process generates the sequence of integers from 2 to 35 and streams them into a pipe.
    \item Each filter stage reads the first integer from its input pipe; by mathematical induction, this integer $p$ is guaranteed to be prime.
    \item The stage prints \texttt{prime p}, creates a new child stage with an outgoing pipe, and filters subsequent integers $n$, forwarding only those where $n \pmod p \neq 0$.
\end{itemize}

\begin{figure}[H]
    \centering
    \small
    \texttt{[Root 2..35] ---> [Stage 2] ---> [Stage 3] ---> [Stage 5] ---> ...}
    \caption{Process pipeline for the concurrent Sieve of Eratosthenes.}
\end{figure}

\subsubsection{Implementation Details in xv6}
In xv6, processes have a strict limit on open file descriptors (\texttt{NOFILE = 16}). Furthermore, recursive child spawning in bare-metal compilers like modern GCC 16 emits infinite-recursion warnings if not carefully managed.

To guarantee zero warnings and deterministic descriptor management, we implemented \texttt{user/primes.c} using an iterative stage handoff loop:

\begin{lstlisting}[style=c]
#include "kernel/types.h"
#include "user/user.h"

void
sieve(int pleft[2])
{
  while(1){
    close(pleft[1]);
    int prime;
    if(read(pleft[0], &prime, sizeof(prime)) != sizeof(prime)){
      close(pleft[0]);
      exit(0);
    }
    printf("prime %d\n", prime);

    int pright[2];
    if(pipe(pright) < 0){
      fprintf(2, "pipe failed\n");
      exit(1);
    }

    int pid = fork();
    if(pid < 0){
      fprintf(2, "fork failed\n");
      exit(1);
    }
    if(pid == 0){
      close(pleft[0]);
      pleft[0] = pright[0];
      pleft[1] = pright[1];
      continue;
    } else {
      close(pright[0]);
      int n;
      while(read(pleft[0], &n, sizeof(n)) == sizeof(n)){
        if(n % prime != 0){
          write(pright[1], &n, sizeof(n));
        }
      }
      close(pleft[0]);
      close(pright[1]);
      wait(0);
      exit(0);
    }
  }
}

int
main(int argc, char *argv[])
{
  int p[2];
  if(pipe(p) < 0){
    fprintf(2, "pipe failed\n");
    exit(1);
  }

  int pid = fork();
  if(pid < 0){
    fprintf(2, "fork failed\n");
    exit(1);
  }
  if(pid == 0){
    sieve(p);
  } else {
    close(p[0]);
    for(int i = 2; i <= 35; i++){
      write(p[1], &i, sizeof(i));
    }
    close(p[1]);
    wait(0);
  }
  exit(0);
}
\end{lstlisting}

\subsubsection{Verification in QEMU}
Running \texttt{primes} inside xv6 produced all prime numbers between 2 and 35:

\begin{lstlisting}[style=bash]
$ primes
prime 2
prime 3
prime 5
prime 7
prime 11
prime 13
prime 17
prime 19
prime 23
prime 29
prime 31
$ 
\end{lstlisting}
Every prime was accurately identified, and all stages exited cleanly after encountering EOF on their respective pipe descriptors.

% ==============================================================================
\section{Conclusion}

In this laboratory work, I evaluated process management and inter-process communication directly inside \textbf{xv6-riscv}, verifying the fundamental principles explored on Linux in Laboratory 7.

Implementing \texttt{user/run.c} demonstrated that the core Unix system calls---\texttt{fork()}, \texttt{exec()}, and \texttt{wait()}---operate under the exact same semantic model in a lightweight educational kernel as they do on a full-scale Linux system. We confirmed that \texttt{fork()} clones the caller's process state, \texttt{exec()} discards the previous memory space to execute the target ELF binary, and \texttt{wait()} synchronizes parent execution with child termination while reaping process state.

Developing \texttt{user/pingpong.c} highlighted the mechanics of pipe-based IPC. Because pipes are strictly unidirectional, two distinct pipes are mandatory to avoid race conditions and self-reads. Furthermore, the blocking nature of \texttt{read()} on empty pipes naturally provides process synchronization without requiring explicit sleep loops or busy-waiting.

Finally, implementing the concurrent prime sieve in \texttt{user/primes.c} demonstrated the power and elegance of Unix pipelines, where complex algorithmic problems decompose into a chain of autonomous filtering processes communicating through kernel FIFOs. This laboratory solidified my understanding of how low-level system call primitives compose to provide robust process isolation and coordination.

\end{document}
